\section{The duality principle}
\label{sec:duality-principle}

The protected determinant is a duality-covariant automorphic section.
The product chamber is a cusp trivialisation of that section.

For
$$
Z=\begin{pmatrix}z_1&z_2\\z_2&z_3\end{pmatrix}\in\Hpl_2,
\qquad
q=e^{2\pi iz_1},\quad r=e^{2\pi iz_2},\quad s=e^{2\pi iz_3},
$$
put
$$
\Gamma_{\mathrm{BPS}}
=\left\{
T(n,l,m)=
\begin{pmatrix}
n&l/2\\ l/2&m
\end{pmatrix}
\ \middle|\ n,l,m\in\Z
\right\}.
$$
The coordinate character is
$$
(T(n,l,m),Z)_{\mathrm{tr}}
:=\operatorname{tr}(T(n,l,m)Z)=nz_1+lz_2+mz_3,
$$
$$
x^{T(n,l,m)}=\exp(2\pi i(T(n,l,m),Z)_{\mathrm{tr}})
=q^nr^ls^m.
$$

\begin{definition}[Cusp charge pairing]
For
$$\gamma=(n,l,m),\qquad \gamma'=(n',l',m'),$$
define
$$
\langle\gamma,\gamma'\rangle_{\mathrm{BPS}}
=2(nm'+n'm)-ll'.
$$
Then
$$
\Delta(\gamma):=\langle\gamma,\gamma\rangle_{\mathrm{BPS}}
=4nm-l^2.
$$
\end{definition}

\begin{definition}[Standard effective chamber]
\label{def:standard-effective-chamber}
The standard cusp is the ordered limit
$$
0<|s|\ll |q|\ll |r|^{-1}\ll 1.
$$
On imaginary parts this is the asymptotic sector
$$
y_3\gg y_1\gg -y_2>0,\qquad y_1y_3-y_2^2>0.
$$
The corresponding positive charge semigroup is
$$
\Gamma_{\mathrm{eff}}
=\{(n,l,m)\mid m>0,\ n\geq0,\ l\in\Z\}
$$
$$
\cup\{(n,l,0)\mid n>0,\ l\in\Z\}
\cup\{(0,l,0)\mid l<0\}.
$$
We write
$$
\gamma>_\infty0\quad\Longleftrightarrow\quad
\gamma\in\Gamma_{\mathrm{eff}}.
$$
\end{definition}

\begin{remark}[Status of the chamber]
The pairing $\langle\ ,\ \rangle_{\mathrm{BPS}}$ is fixed by the
duality lattice below.  The order $>_\infty$ is a convention: it is the
Fock polarisation at the standard zero-dimensional cusp.  A duality
element sends it to the product chamber at the transported cusp; it does
not preserve the displayed lexicographic inequalities inside the same
coordinate chart.
\end{remark}

The Lorentzian degree lattice used by the denominator algebra is
$$
\La^{2,1}_{II}
=2\Z f_2\oplus \Z f_3\oplus 2\Z f_{-2},
$$
with
$$
(f_2,f_{-2})=-1,\qquad (f_3,f_3)=2.
$$
Define
$$
\alpha:\Gamma_{\mathrm{BPS}}\longrightarrow\La^{2,1}_{II},
\qquad
\alpha(n,l,m)=2nf_2-lf_3+2mf_{-2}.
$$
Then
$$
(\alpha(\gamma),\alpha(\gamma'))
=-2\langle\gamma,\gamma'\rangle_{\mathrm{BPS}},
\qquad
(\alpha(\gamma),\alpha(\gamma))=-2(4nm-l^2).
$$
For
$$
z=z_1f_{-2}+z_2f_3+z_3f_2
$$
one has
$$
(\alpha(n,l,m),z)=-2(nz_1+lz_2+mz_3),
$$
hence
$$
q^nr^ls^m
=\exp(2\pi i(nz_1+lz_2+mz_3))
=\exp(-\pi i(\alpha(n,l,m),z)).
$$
Thus the physical discriminant $4nm-l^2$ is the negative half of the
Lorentzian square of the BKM degree.

\begin{definition}[Duality lattice]
Let
$$
\La^4=\Z e_1\oplus\Z e_2\oplus\Z e_3\oplus\Z e_4,
\qquad
\operatorname{vol}=e_1\wedge e_2\wedge e_3\wedge e_4.
$$
On $\wedge^2\La^4$ put
$$
X\wedge Y=(X,Y)\operatorname{vol}.
$$
Then $(\wedge^2\La^4,(\ ,\ ))\simeq \mathrm{II}_{3,3}$.
Fix
$$
q_0=e_1\wedge e_3+e_2\wedge e_4
$$
and
$$
\La_{\mathrm{du}}
:=q_0^\perp\subset\wedge^2\La^4.
$$
The basis
$$
f_1=e_1\wedge e_2,\quad
f_2=e_2\wedge e_3,\quad
f_3=e_1\wedge e_3-e_2\wedge e_4,\quad
f_{-2}=e_4\wedge e_1,\quad
f_{-1}=e_4\wedge e_3
$$
gives
$$
\La_{\mathrm{du}}\simeq\La^{3,2},\qquad
(f_1,f_{-1})=(f_2,f_{-2})=-1,\qquad (f_3,f_3)=2.
$$
\end{definition}

The Siegel point $Z\in\Hpl_2$ is the isotropic line generated by
$$
\mathcal Z(Z)
=(z_2^2-z_1z_3)f_1+z_3f_2+z_2f_3+z_1f_{-2}+f_{-1}.
$$
Then
$$
(\mathcal Z(Z),\mathcal Z(Z))=0,
$$
and
$$
(\mathcal Z(Z),\overline{\mathcal Z(Z)})<0
\quad\Longleftrightarrow\quad
\Imag Z>0
\quad\Longleftrightarrow\quad
y_1y_3-y_2^2>0,\ y_1>0.
$$
For $\gamma=(n,l,m)$,
$$
(\alpha(\gamma),\mathcal Z(Z))
=-2(nz_1+lz_2+mz_3),
$$
so the cusp charges pair with the full period line by the same
character:
$$
x^\gamma=\exp(-\pi i(\alpha(\gamma),\mathcal Z(Z))).
$$

\begin{theorem}[Exterior-square realisation of the duality group]
\label{thm:duality-wedge-square}
Let
$$
B_{q_0}(x,y)\operatorname{vol}:=-x\wedge y\wedge q_0.
$$
Then
$$
B_{q_0}(e_1,e_3)=1,\qquad B_{q_0}(e_2,e_4)=1,
$$
and
$$
\{g\in\SL(\La^4)\mid (\wedge^2g)q_0=q_0\}
=\Sp(\La^4,B_{q_0})\simeq\Sp_4(\Z).
$$
The exterior-square representation gives
$$
1\longrightarrow\{\pm\Id_4\}
\longrightarrow\Sp_4(\Z)
\xrightarrow{\ \wedge^2\ }
\SO_+(\La^{3,2})
\longrightarrow 1,
$$
equivalently
$$
\PSp_4(\Z)\simeq\SO_+(\La^{3,2})
\simeq\Orth(\La^{3,2})_+/\{\pm\Id_5\}.
$$
For
$$
g=\begin{pmatrix}A&B\\ C&D\end{pmatrix}\in\Sp_4(\Z),
\qquad
g\cdot Z=(AZ+B)(CZ+D)^{-1},
$$
one has the projective identity
$$
(\wedge^2g)\mathcal Z(Z)
=\det(CZ+D)\,\mathcal Z(g\cdot Z).
$$
\end{theorem}

\begin{proof}
For $g\in\SL(\La^4)$,
$$
(\wedge^2gX)\wedge(\wedge^2gY)=X\wedge Y,
$$
so $\wedge^2g$ preserves the Pfaffian form on $\wedge^2\La^4$.
The equation $(\wedge^2g)q_0=q_0$ is equivalent to
$$
B_{q_0}(gx,gy)=B_{q_0}(x,y),
$$
which is the defining equation of $\Sp(\La^4,B_{q_0})$.

Write
$$
u_1(Z)=z_1e_1+z_2e_2+e_3,\qquad
u_2(Z)=z_2e_1+z_3e_2+e_4.
$$
Then
$$
\mathcal Z(Z)=-u_1(Z)\wedge u_2(Z).
$$
Since
$$
g\begin{pmatrix}Z\\ \Id_2\end{pmatrix}
=
\begin{pmatrix}AZ+B\\ CZ+D\end{pmatrix}
=
\begin{pmatrix}g\cdot Z\\ \Id_2\end{pmatrix}(CZ+D),
$$
taking the wedge of the two columns gives
$$
(\wedge^2g)\mathcal Z(Z)
=\det(CZ+D)\mathcal Z(g\cdot Z).
$$
The kernel on $\La^{3,2}=q_0^\perp$ is $\{\pm\Id_4\}$, and the
exceptional isomorphism $\PSp_4(\R)\simeq\SO_{\R}(3,2)_+$ identifies
the integral image with $\SO_+(\La^{3,2})$.
\end{proof}

\begin{definition}[Automorphic duality datum]
The geometric duality group is
$$
G_{\mathrm{du}}=\PSp_4(\Z)\simeq\SO_+(\La^{3,2}).
$$
The chiral determinant is not an invariant scalar function on $\Hpl_2$.
It is a section of the automorphic line
$$
\mathcal L^5\otimes\nu_{\Delta_5}
$$
with transformation law
$$
\Delta_5(g\cdot Z)
=\nu_{\Delta_5}(g)\det(CZ+D)^5\Delta_5(Z),
\qquad
g=\begin{pmatrix}A&B\\ C&D\end{pmatrix}\in\Sp_4(\Z).
$$
The scalar square is a section of $\mathcal L^{10}$:
$$
\Delta_5^2(g\cdot Z)=\det(CZ+D)^{10}\Delta_5^2(Z).
$$
Consequently
$$
Z^X_{\square}(g\cdot Z)
=\det(CZ+D)^{-10}Z^X_{\square}(Z)
$$
for
$$
Z^X_{\square}=C_{\square}\Delta_5^{-2}.
$$
\end{definition}

The duality-invariant wall arrangement is the divisor, not a chosen
equation in one cusp.  Gritsenko--Clery \cite{GC}, $N=1$:
$$
\mathcal M_{\mathrm{massless}}
=\operatorname{supp}\operatorname{div}(\Delta_5)
=\Sp_4(\Z)\cdot\mathcal H_{\mathrm{diag}},
$$
$$
\mathcal H_{\mathrm{diag}}
=\left\{
\begin{pmatrix}a&0\\0&b\end{pmatrix}
\ \middle|\ a,b\in\mathbb H_1
\right\},
\qquad
\operatorname{mult}_{\mathcal H_{\mathrm{diag}}}(\Delta_5)=1.
$$

\begin{proposition}[Duality-covariant product in the standard chamber]
In the standard cusp chamber,
$$
\mathcal D_X(Z)
=64q^{1/2}r^{1/2}s^{1/2}
\prod_{\gamma=(n,l,m)>_\infty0}
(1-q^nr^ls^m)^{f(nm,l)}.
$$
Equivalently,
$$
\mathcal D_X(Z)
=64\exp(-\pi i(\rho,z))
\prod_{\gamma>_\infty0}
\left(1-\exp(-\pi i(\alpha(\gamma),z))\right)^{f(nm,l)}.
$$
The equality
$$
\mathcal D_X(Z)=\Delta_5(Z)
$$
upgrades this chamber product to the automorphic section
$$
\mathcal D_X(g\cdot Z)
=\nu_{\Delta_5}(g)\det(CZ+D)^5\mathcal D_X(Z).
$$
\end{proposition}

\begin{proof}
The first formula is the Borcherds--Gritsenko--Nikulin product for the
weak Jacobi form $\phi_{0,1}$ in the standard cusp \cite{GN,B}.  The second follows
from
$$
q^nr^ls^m=\exp(-\pi i(\alpha(n,l,m),z)),
\qquad
q^{1/2}r^{1/2}s^{1/2}=\exp(-\pi i(\rho,z)).
$$
The last displayed equation is Maass automorphy of $\Delta_5$
\cite{MAASS:1}.
\end{proof}

\begin{remark}[Theorem, convention, assumption]
The following are the precise statuses used here.
\begin{itemize}
\item Theorem: the charge pairing satisfies
$$
\langle\gamma,\gamma'\rangle_{\mathrm{BPS}}
=-\frac12(\alpha(\gamma),\alpha(\gamma')).
$$
\item Theorem: the duality group acting on the period domain is
$$
\PSp_4(\Z)\simeq\SO_+(\La^{3,2})
$$
via $\wedge^2$.
\item Theorem: $\Delta_5$ is a weight-$5$ Siegel modular form with
Maass character $\nu_{\Delta_5}$ \cite{MAASS:1}, divisor
$$
\Sp_4(\Z)\cdot\mathcal H_{\mathrm{diag}},
$$
and product expansion in the standard effective chamber
\cite{GN,B,GC}.
\item Convention: the chamber $\Gamma_{\mathrm{eff}}$ is the standard
cusp/Fock polarisation
$$
0<|s|\ll |q|\ll |r|^{-1}\ll1.
$$
\item Convention: the sign of $\Delta_5$ is fixed by
$$
[q^{1/2}r^{1/2}s^{1/2}]\Delta_5=64.
$$
\item Assumption: the Donaldson--Thomas scalar-square interpretation is
only the assertion
$$
Z^X_{\mathrm{DT}}=C_{\mathrm{DT}}\Delta_5^{-2}.
$$
The duality lattice, the chamber product, and the automorphic covariance
do not depend on this assumption.
\end{itemize}
\end{remark}
